\documentclass[11pt]{article}
\usepackage{fontspec}
\usepackage[french]{babel}
\usepackage{amsmath}
\usepackage{amssymb}
\usepackage{amsthm}
\usepackage{mathrsfs}
\usepackage{tikz-cd}
\usepackage{enumitem}
\usetikzlibrary{babel}

\newcommand{\printedpage}[1]{\marginpar{\footnotesize p.~#1}}
\newcommand{\msp}[1]{\marginpar{\footnotesize [#1]}}
\newcommand{\cercle}[1]{\tikz[baseline=(c.base)]\node[draw,circle,inner sep=1pt](c){#1};}
\newcommand{\defeq}{\overset{\text{déf}}{=}}
\newcommand{\upi}{\underline{\pi}}
\newcommand{\Cat}{\mathrm{Cat}}
\newcommand{\Ob}{\mathrm{Ob}\,}

\begin{document}

\printedpage{63}
% head: A. Grothendieck : Dérivateurs, Ch. XII, Édition des Universités de Montpellier II et Paris VII
et on aura
\[ U'/t' \simeq U/H'. \]
La relation entre $H'$ et $H$ a été explicitée au §1, comme
\[ H' = p_!H, \qquad \text{où } p_! : U^{\wedge} \longrightarrow U'^{\wedge} \]
est l'adjoint à gauche de $p^* : U'^{\wedge} \longrightarrow U^{\wedge}$, donné par
\[ p^*F' = F' \circ p. \]
On a
\[ p_!(H)(u') \simeq \pi_0(u'\backslash U/H). \]

\textbf{Définition 1.} \cercle{a} Soit $U$ un objet de Cat, $H \in \Ob U^{\wedge}$ un préfaisceau sur $U$. On dit que $H$ est \textit{totalement $W$-asphérique} si pour toute flèche $p : U \longrightarrow U'$ de Cat, $p_!(H)$ dans $U'^{\wedge}$ est un objet $W$-asphérique, i.e. $U'/p_!H$ est $W$-asphérique. \cercle{b} On dit qu'un objet $V$ de \msp{page 89} Cat est \textit{totalement $W$-asphérique,} si le foncteur final $e_V$ sur $V$ ($v \longmapsto \{\emptyset\}$) est totalement $W$-asphérique, i.e. si pour toute flèche $f : V \longrightarrow V'$ dans Cat, le préfaisceau
\[ \pi_0(f) \defeq f_!(e_V) \defeq (v' \longmapsto \pi_0(v'\backslash V)) \]
sur $V'$ est $W$-asphérique, i.e. si la catégorie
\[ \upi_0(f) \ \text{ ou } \ \upi_0(V/V') \defeq V'/\pi_0(f) \]
est $W$-asphérique. \texttt{[Cette notion est différente de la notion de totale asphéricité de Pursuing Stacks.]}

Notons que la définition des foncteurs $p_!$ (ou $f_!$) comme adjoints, et leur `calcul' en termes de $\pi_0$ de certaines catégories, ne dépend pas du localiseur $W$. Donnons quelques précisions à ce sujet. On a un morphisme d'adjonction dans $V^{\wedge}$
\[ f^*f_!(F) \longleftarrow F, \]
d'où
\[ V/F \longrightarrow V/f^*f_!(F) \simeq V \times_{V'} (V'/f_!(F)), \tag*{$(*)$} \]
le dernier isomorphisme provenant du fait que l'on a commutativité, à isomorphisme canonique près, dans le diagramme de foncteurs
\[
\begin{tikzcd}[column sep=14em, row sep=large]
V'^{\wedge} \arrow[r, "f^*=(F'\longmapsto F'\circ f)"] \arrow[d, "F'\longmapsto V'/F'"'] & V^{\wedge} \arrow[d, "F\longmapsto F/V"] \\
\Cat/V' \arrow[r, "f^*_{\Cat}=(W'\longmapsto W'\times_{V'}V)"] & \Cat/V,
\end{tikzcd}
\]
% foot: Version 7 décembre 2002
\printedpage{64}
% head: A. Grothendieck : Dérivateurs, Ch. XII, Édition des Universités de Montpellier II et Paris VII
\msp{page 90}
en d'autres termes, l'opérateur $f^*$ sur les préfaisceaux s'interprète, en termes géométriques, comme l'opération $W' \longmapsto W' \times_{V'} V$ de changement de base par $V \longrightarrow V'$, sur les catégories au dessus de $V'$ qui sont fibrées à fibres discrètes (propriété qui se conserve par changement de base).

Ceci dit, la donnée d'une flèche $(*)$ équivaut, par la propriété universelle du changement de base, à celle d'une flèche (canonique en l'occurrence)
\[ f/F : V/F \longrightarrow V'/f_!(F) \]
rendant commutatif le carré
\[
\begin{tikzcd}[column sep=large, row sep=huge]
V/F \arrow[r, "\text{oubli de }F", "\substack{\text{ou foncteur}\\ \text{source}}"'] \arrow[d, "f/F"'] & V \arrow[d, "f"] \\
V'/(F' = f_!(F)) \arrow[r, "\text{oubli de }F'"] & V'.
\end{tikzcd}
\tag*{$(*)$}
\]
Dans le cas qui nous intéresse surtout ici, $F = e_V$, $V/F \longrightarrow V$ est un isomorphisme, et on obtient une factorisation canonique de foncteurs \msp{page 91}
\[
\begin{tikzcd}[row sep=large]
 & V \arrow[dl, "\tilde{f}=f/e_V"'] \arrow[dd, "f"] \\
\begin{array}{r} \upi_0(V/V') \\ = V'/f_!(e_V) \end{array} \arrow[dr, "\text{can.}"'] & \\
 & V'.
\end{tikzcd}
\tag*{$(**)$}
\]
Explicitons le $V'$-morphisme $\tilde{f}$ sur les fibres,
\[ \tilde{f}_{v'} : V_{v'} \longrightarrow \upi_0(V/V')_{v'} \defeq \pi_0(v'\backslash V), \tag*{$(***)$} \]
où dans le dernier membre, on a identifié l'ensemble $\pi_0(v'\backslash V)$ à la catégorie discrète qu'il définit. Pour expliciter cette flèche $(***)$, notons qu'on a une inclusion canonique
\[ V_{v'} \longrightarrow v'\backslash V \]
et pour toute catégorie $W$, un morphisme canonique (d'adjonction) $W \longrightarrow \pi_0(W)$, d'où un composé
\[ V_{v'} \longrightarrow v'\backslash V \longrightarrow \pi_0(v'\backslash V), \]
qui n'est autre que $\tilde{f}_{v'}$ dans $(***)$.

Il doit être facile de montrer que la factorisation $(**)$ de $f$ est caractérisée, à isomorphisme unique près, en tant que
% foot: Version 7 décembre 2002
\printedpage{65}
% head: A. Grothendieck : Dérivateurs, Ch. XII, Édition des Universités de Montpellier II et Paris VII
\msp{page 92}
factorisation dans $\Cat/V'$
\[
\begin{tikzcd}[column sep=small]
 & \tilde{V} \arrow[dr, "\pi"] & \\
V \arrow[ur, "\tilde{f}"] \arrow[rr, "f"'] & & V',
\end{tikzcd}
\]
par les propriétés suivantes :
\begin{enumerate}[label={}, leftmargin=*]
\item[a)] $\pi$ est fibrée à fibres discrètes, i.e. un isomorphisme local (i.e. $\tilde{V}$ est isomorphe dans $\Cat/V'$ à une catégorie $V/F'$, avec $F'$ dans $V'^{\wedge}$). % [sic] V/F' -> V'/F', une catégorie au-dessus de V' définie par F' dans V'^∧ est de la forme V'/F'
\item[b)] $\tilde{f}_!(e_V) = e_{\tilde{V}}$, i.e. pout tout $\tilde{v} \in \Ob \tilde{V}$, la catégorie $\tilde{v}\backslash V$ est 0-connexe : les colocalisés de $V$ au dessus de $\tilde{V}$ sont 0-connexes. % [sic] pout -> pour, faute de frappe
\end{enumerate}

Notons ici, en passant, que pour \textit{tout} morphisme $\pi : W' \longrightarrow W$ dans Cat qui soit fibré à fibres discrètes, i.e. tel que $W'$ soit isomorphe dans $\Cat/W$ à une catégorie de la forme $W/F$ ($F \in \Ob W^{\wedge}$), on a
\[ \pi_!(e_{W'}) = F, \qquad \text{i.e. } W' \underbrace{\ \simeq\ }_{\text{dans } \Cat/W} W/\pi_!(e_{W'}). \tag*{$(*)$} \]
Il faut vérifier que le foncteur
\[ w \longmapsto \pi_0(\underbrace{w\backslash(W/F)}_{\defeq w\backslash W/F}) \]
est canoniquement isomorphe à $F$. Mais l'inclusion
\[ \underbrace{W'_w}_{= F(w) \text{ (catégorie discrète)}} \qquad \longrightarrow w\backslash W' \]
\msp{page 93}
définit bien comme tantôt
\[ F(w) = W'_w \longrightarrow \pi_0(w\backslash W'), \]
il faut prouver que c'est bijectif. Or, $W'$ étant $W_\infty$-lisse sur $W$, on sait que l'inclusion $W'_w \hookrightarrow w\backslash W'$ est dans $W_\infty$, a fortiori induit un isomorphisme pour les $\pi_0$.

Ceci posé, si on a une factorisation de $f : V \longrightarrow V'$ en $f = \pi\tilde{f}$ ayant les propriétés a) b) de la page précédente, on voit déjà que l'on a
\[ f_!(e_V) \simeq \pi_!(\, \underbrace{\tilde{f}_!(e_V)}_{\simeq e_{\tilde{V}} \text{ par b)}} \,) \simeq \pi_!(e_{\tilde{V}}), \]
ce qui déjà détermine $\tilde{V}$ à isomorphisme canonique près (grâce à a) comme
\[ \tilde{V} \simeq V'/f_!(e_V) \defeq \upi_0(V'/V), \] % [sic] \underline{\pi}_0(V'/V) -> \underline{\pi}_0(V/V'), notation employée partout ailleurs (Définition 1, (**))
% foot: Version 7 décembre 2002
\printedpage{66}
% head: A. Grothendieck : Dérivateurs, Ch. XII, Édition des Universités de Montpellier II et Paris VII
et il ne doit pas être bien dur de se convaincre que $\tilde{f}$ n'est autre que la flèche canonique de la page 91.

\textbf{Proposition 8.} \textit{Pour qu'un objet $H$ de $U^{\wedge}$ soit totalement $W$-asphérique, il suffit que l'objet $U/H$ de} Cat \textit{soit totalement $W$-asphérique.}

\msp{page 94}
Supposons en effet cette dernière condition satisfaite et prouvons que $H$ est totalement $W$-asphérique, i.e. que pour tout $p : U \longrightarrow U'$, $U'/p_!(H)$ soit $W$-asphérique.
\[
\begin{tikzcd}[row sep=large]
V \defeq U/H \arrow[r, "\alpha"] \arrow[dr, "f"'] & U \arrow[d, "p"] \\
 & U'
\end{tikzcd}
\]
Soit $f$ le foncteur composé du diagramme ci-contre \texttt{[ci-dessus]}, on a $H \simeq \alpha_!(e_V)$ (par $(*)$ p. 92), d'où
\[ p_!(H) \simeq p_!\alpha_!(e_V) \simeq (p\alpha)_!(e_V) \simeq f_!(e_V), \]
qui est $W$-asphérique par hypothèse sur $V = U/H$.

La réciproque de \texttt{[la]} prop. 8 est-elle valable, i.e. si $H$ est totalement $W$-asphérique, en est-il de même de $V = U/H$? La difficulté provient du fait que cela signifie que pour \textit{toute} flèche $f : V \longrightarrow V'$ dans Cat de source $V$, on ait $f_!(e_V)$ $W$-asphérique, alors que l'hypothèse équivaut à le postuler seulement pour les flèches qui se factorisent par $U$. Notons que si on ne fait pas cette hypothèse sur $f$, mais qu'on suppose au lieu de cela que $f$ s'insère dans un carré commutatif \msp{page 95}
\[
\begin{tikzcd}[row sep=large]
V \arrow[r, "\alpha"] \arrow[d, "f"'] & U \arrow[d, "p"] \\
V' \arrow[r, "\alpha'"] & U',
\end{tikzcd}
\]
avec $\alpha'$ lui aussi un isomorphisme local, alors on peut conclure que $\upi_0(f) = \upi_0(V/V')$ est bien asphérique. En effet, on a un isomorphisme canonique dans $\Cat/U'$,
\[ \upi_0(\alpha'f) \simeq \upi_0(f), \qquad \text{i.e. } \ \upi_0(V/V') \simeq \upi_0(V/U'), \]
du fait p. ex. que la factorisation canonique de $f$
% foot: Version 7 décembre 2002
\printedpage{67}
% head: A. Grothendieck : Dérivateurs, Ch. XII, Édition des Universités de Montpellier II et Paris VII
\[
\begin{tikzcd}[row sep=normal]
 & V \arrow[dl, "\tilde{f}"'] \arrow[dd, "f"] & \\
\upi_0(V/V') \arrow[dr, "\pi"'] & & \\
 & V' \arrow[dr, dotted, "\alpha'"'] & \\
 & & U'
\end{tikzcd}
\]
fournit une factorisation du même type pour $\alpha'f$, en
\[ \alpha'f = (\alpha'\pi) \circ \tilde{f}. \]
Ceci dit, on peut appliquer l'hypothèse à $p$ et on trouve bien la $W$-asphéricité de $\upi_0(f) \simeq \upi_0(\alpha'f) = \upi_0(p\alpha)$.

\msp{page 96}
L'idée naturelle, pour une réciproque à \texttt{[la]} prop. 8, serait de prendre pour $U'$ la somme amalgamée de $U$ et $V'$ sur $V$, ce qui marchera pourvu qu'on sache que $\alpha'$ est bel et bien un isomorphisme local tout comme $\alpha$. C'est bien le cas p. ex. si $V \longrightarrow U$ est une immersion ouverte, i.e. $H$ un sous-préfaisceau du préfaisceau final. Mais je crains que ce ne soit pas vrai en général - ça a l'air un peu trop beau pour être vrai! C'est en effet déjà faux dans le cas
\[
\begin{tikzcd}[row sep=huge]
V = E_G \arrow[r, "\alpha\ \text{can}"] \arrow[d, "f"'] & B_G = U \\
V' = B_{G'}, &
\end{tikzcd}
\]
dans le cas p. ex. où $G = \{\pm 1\}$, groupe à deux éléments, et où l'élément $\alpha$ de $G'$, image par $f$ de l'une des deux flèches non identiques de $E_G$ est $\neq 1$ \texttt{[plutôt de carré} $\neq 1$\texttt{]}. Si $H'$ est le sous-groupe invariant engendré par $\alpha$ \texttt{[plutôt par} $\alpha^2$\texttt{]} dans $G'$, la somme amalgamée s'identifie à $B_{G'/H'}$, or $B_{G'} \longrightarrow B_{G'/H'}$ n'est \textit{pas} un isomorphisme local (car $(B_{G'})/t_0 \longrightarrow (B_{G'/H'})/t'_0$ n'est pas bijectif sur les objets, s'identifiant à $G' \overset{q}{\longrightarrow} G'/H' \ \ldots$).

Mais d'autre part la condition de totale $W$-asphéricité sur un objet $H$ de $U^{\wedge}$ paraît déjà à tel point exigeante, que je me vois mal con\msp{page 97}struire un exemple où elle est satisfaite, sans que $U/H$ elle-même ne soit totalement $W$-asphérique.

L'existence d'objets totalement $W$-asphériques dans Cat n'est pas tellement évidente a priori. Qu'il y en ait un bon paquet est assuré par la
% foot: Version 7 décembre 2002

\end{document}
